\documentclass[11pt]{article}
\usepackage[T1]{fontenc}
\usepackage{lmodern,amsmath,amssymb,amsthm}
\usepackage[a4paper,margin=27mm]{geometry}
\usepackage[hidelinks]{hyperref}
\newtheorem{proposition}{Proposition}
\newtheorem{corollary}[proposition]{Corollary}
\newcommand{\rev}{\operatorname{rev}}
\title{Which part of contour word solving is intrinsically difficult?}
\author{CPLX-002 --- LLM-generated research note and proof drafts}
\date{21 October 2026}
\begin{document}
\maketitle

\noindent\textbf{Status.} The elementary reductions below are proof drafts,
not independently accepted results or enabled solver replacements. They classify
specific finite algebraic tasks, not the existence of disk tilings. Their
elementary periodic construction is not claimed as a new literature theorem.
The interval-orbit algorithm cited below is published; its application to our
signed position constraints is derived here. No end-to-end speedup is claimed.

\section{Extracted input, and a crucial distinction}

At fixed finite incidence and weak contour order, the compiler in
\texttt{words/compile.py} creates one variable $X_i$ per atomic prototype
interval. The prototype is the cyclic sequence $X_0\cdots X_{p-1}$.
An interface equates two oriented cyclic factors, possibly on separate direct
and mirror sheets. Each is a proper factor without repeating an atomic
variable on that side. Across all equations a variable may appear many times;
this is not a guarantee of a globally quadratic system. Reflections introduce
variables $M_i$, and optional intrinsic symmetry introduces a sector variable
and a power equation. None is a \emph{constant letter}.

An occurrence $X^{-1}$ means an anti-involution $\mu$: reverse the sequence
and exchange every letter with its opposite. There is \emph{no free-group
cancellation} $a\bar a=1$. For instance, $a\bar a$ is a nonempty word fixed
by $\mu$. This is also how \texttt{words/algebra.py} and the palindrome branch
operate. In particular, the input is not an arbitrary word equation with
prescribed constant strings or arbitrary regular constraints.

The compiler stores outer-arc words as metadata; their curvature, endpoint and
angle requirements are not constant letters in the bare word equation.
Contact V2 also propagates explicit corners and metric constraints before the
leaf solver. Neither that decorated partial state nor the full geometric
problem is classified by the bare-word statements below.

\section{Constant-free existence is linear feasibility}

Let $E$ be a finite system of equations between words in variables
$X_1,\ldots,X_p$ and their $\mu$-images, with no constant letters. All
variable images must be nonempty finite words in an alphabet containing at
least one opposite pair $a,\bar a$, $\bar a\ne a$.
For each equation make one integer row by counting variable occurrences on
the left minus the right, ignoring inverse flags; let $A$ be the resulting
matrix.

\begin{proposition}[Bare existence]\label{prop:bare}
The following are equivalent:
\begin{enumerate}
\item $E$ has a non-erasing solution over a free monoid with this anti-involution.
\item $A n=0$ has a strictly positive integer solution.
\item $A x=0$ has a strictly positive real solution.
\end{enumerate}
This bare decision problem is in polynomial time in the explicit equation
input size. A witness may require exponentially many letters when expanded;
a witness encoded by binary exponents has polynomial-size description.
\end{proposition}

\begin{proof}
Word lengths prove (1)$\Rightarrow$(2), and (2)$\Rightarrow$(3) is immediate.
For (3)$\Rightarrow$(2), scale a positive real solution so every coordinate
is at least one. The nonempty rational polyhedron $A x=0$, $x_i\geq1$
has a rational point. Clearing denominators gives positive integers $n_i$.

Put $w=a\bar a$, so $\mu(w)=w$, and assign $X_i=w^{n_i}$. Both $X_i$
and $\mu(X_i)$ become the same power of $w$. Each equation becomes equality
of powers whose exponents differ by one row of $A n=0$. This proves (1).

Rational linear feasibility is polynomial-time decidable. Standard rational
polyhedron bounds give a feasible rational vector of polynomial encoding
length; clearing its denominators still has polynomial encoding length.
The assignment $X_i=(a\bar a)^{n_i}$ therefore gives a succinct witness.
This is not a claim that the repository's simplex implementation has a
polynomial worst-case pivot count.
\end{proof}

Additional homogeneous rational constraints on the lengths can be included in
the same argument when they are invariant under positive scaling. Mirror-pair
length equalities are an example. Fixed constants, prescribed noncommuting
letter types, exact prescribed odd word lengths, reduced words, fixed alphabet
restrictions excluding $a\bar a$, and point/curve decorations are not silently
included. In particular $X=\mu(X)$ has a nonempty solution but has none of
\emph{exactly one letter} over the stated fixed-point-free alphabet.

As a succinctness control, the constant-free system
$X_i=X_{i+1}X_{i+2}$ ($0\leq i\leq n-2$), $X_{n-1}=X_n$, has minimum
positive lengths growing as Fibonacci numbers. Its size is linear in the
number of variable occurrences, up to identifier encoding. All equation sides
are proper factors of the same cyclic order $X_0\cdots X_n$ when $n\geq3$.
This shows exponential \emph{expanded output} even in that syntactic factor
class; it is not a proof that this family comes from valid embedded tile maps,
nor a hardness result for tiling or for succinct existence.

\section{Why the cheap witness does not solve the disk problem}

\begin{proposition}[Loss of a signed additive observable]\label{prop:charge}
Let $\chi$ be any additive real-valued observable on terminal words with
$\chi(\mu U)=-\chi(U)$. The witness in Proposition \ref{prop:bare} assigns
zero $\chi$ to every variable and every compiled word.
It cannot satisfy a requirement $\chi(U)>0$ on an outer factor.
\end{proposition}

\begin{proof}
$\chi(a\bar a)=\chi(a)-\chi(a)=0$. Additivity proves the assertion.
\end{proof}

An oriented circular-radius measure in the finite decorated-curve model is
such an observable: a nondegenerate positively traversed exterior circle arc
has positive measure, and reversing traversal negates it. This is a typed
geometric constraint, not a bare constant-free equation. We do not identify
word length with physical arclength, nor assert that all adjacent opposite
curve-template occurrences geometrically retrace one another.

Even linear observables do not capture every consequence. The classical
commutation fact $XY=YX$ implies $X=u^s$, $Y=u^t$ for a nonempty $u$ and
positive integers $s,t$. A short proof removes a shorter common prefix and
inducts on $|X|+|Y|$. For any positive additive metric $\ell$ and additive
observable $\chi$, this forces
\[
 \chi(X)\ell(Y)=\chi(Y)\ell(X).
\]
Length balance and the signed additive row of $XY=YX$ alone do not express
that consequence: both rows are identities. For example equal positive
metric lengths and different prescribed $\chi$ are incompatible with
commutation. This is an explanatory algebraic example, not an asserted
compiler-generated disk obstruction or a new commutation theorem.

\section{Fixed word lengths admit a compressed positional test}

Now prescribe positive \emph{integer word lengths} $L_i$ in binary and let
$N=\sum_i L_i$. These are not unknown physical arclengths. Store the letter
positions of each variable in one contiguous block of $[1,N]$.
Each aligned part of an equation gives a bijection between two integer
subintervals, order-preserving or order-reversing, and a sign
$\epsilon\in\{0,1\}$ meaning that the corresponding letters agree or are
opposites. Merge the factor boundaries on the two equation sides; this creates
only a linear number of such pairings in the explicit equation size, without
expanding the $L_i$ positions. Unequal total side lengths reject immediately.

Let $c$ count orbits of these pairings with their signs forgotten. Form a
double cover with positions $(p,b)$, $1\leq p\leq N$, $b\in\{0,1\}$,
and lift every relation $p\sim q$ of sign $\epsilon$ to
\[
 (p,b)\sim(q,b\mathbin{\oplus}\epsilon),\qquad b=0,1.
\]
Lay the sheets in successive blocks of $[1,2N]$, giving another collection
of integer interval pairings; let $\widetilde c$ be its orbit count.

\begin{proposition}[Signed double-cover criterion]\label{prop:cover}
The prescribed-length constant-free system is satisfiable if and only if
\[
 \widetilde c=2c.
\]
Using the published Agol--Hass--Thurston interval orbit-counting theorem
\cite[Theorem 12]{AHT}, the test takes polynomial time in the explicit
equation size and $\log N$.
\end{proposition}

\begin{proof}
For each unsigned orbit the double cover has two components unless a closed
path has odd total sign. Such an odd path joins $(p,0)$ to $(p,1)$ and makes
the cover connected over that orbit. It would require one letter to equal its
opposite, forbidden by the alphabet involution. Thus $\widetilde c=2c$
holds exactly when every orbit is sign-consistent. In that case choose $a$
on one lifted component of each orbit and $\bar a$ on its mate; this assigns
consistent letters and solves all equations. Necessity is the same argument
applied to any word solution. Both orbit counts are AHT instances of polynomial
description size, proving the bound.
\end{proof}

This tests consistency, not preservation of every useful distinct curve class,
geometric coordinates, corner angles, circular signs or unknown-length choices.
The number of orbits alone is insufficient input for the geometry engine.
Recovering a geometry-compatible succinct representation remains an obligation;
the theorem does not license guessing one length vector and declaring all
others exhausted. No AHT implementation or practical speedup is supplied here.

\section{Consequences for the research choice}

The current word layer is not justified as a necessary NP-hard bare feasibility
test merely because it traverses many Nielsen branches. Its main role is to
produce factor-identification families for the decorated/geometric stages.
Those families cannot be replaced by the single periodic witness without
discarding possible geometric solutions. The useful next question is whether
the identifications and their circle/corner constraints can be propagated in
a succinct interval representation before enumerating full word families.

General word-equation complexity \cite{Jez,DJP} does not classify this decorated
subclass. The bare result above also does not classify it. A fixed graph leaves
incidence choices, refinement size and metric/geometry choices unresolved.
The next iteration must target that remaining problem, not claim that an LP
now solves monohedral disk tilings.

\begin{thebibliography}{3}
\bibitem{Jez} A. Je\.{z}, \emph{Word equations in linear space},
\href{https://arxiv.org/abs/1702.00736}{arXiv:1702.00736v2}.
\bibitem{DJP} V. Diekert, A. Je\.{z}, W. Plandowski,
\emph{Finding All Solutions of Equations in Free Groups and Monoids with Involution},
\href{https://arxiv.org/abs/1405.5133}{arXiv:1405.5133v2}.
\bibitem{AHT} I. Agol, J. Hass, W. P. Thurston,
\emph{The Computational Complexity of Knot Genus and Spanning Area},
\href{https://arxiv.org/html/math/0205057v2}{arXiv:math/0205057v2}, Section 4,
Theorem 12. Only interval orbit counting is invoked, not knot-genus hardness.
\end{thebibliography}
\end{document}
